\section*{The Decay Hierarchy: The Geometry of Time}

Why are the Electron and Proton infinitely stable, while the Muon decays in microseconds and the Top Quark in yoctoseconds? We resolve this hierarchy by introducing the \textbf{Persistence Margin (PM)} as the geometric resolution floor of the vacuum.

\subsection*{The Stability Gates}

A particle is stable if it satisfies \textbf{either} of two distinct persistence criteria:

\begin{enumerate}
    \item \textbf{Geometric Lock (The Invisibility Cloak):} If $R_{rel} \le PM$, the geometric flaw is below the resolution threshold of the vacuum. The vacuum cannot "see" the error to punish it.
    \begin{itemize}
        \item \textit{Mechanism:} Resolution Limit.
        \item \textit{Archetype:} Electron ($R=0$).
    \end{itemize}
    
    \item \textbf{Topological Ground State (The Energy Floor):} If the particle is the lightest possible configuration with a conserved topological quantum number (e.g., Baryon Number), it cannot decay, regardless of its geometric residual. Decay requires a lower-energy state accessible via a continuous transformation. If no such state exists, the residual represents \textbf{Internal Stress} rather than \textbf{Decay Width}.
    \begin{itemize}
        \item \textit{Mechanism:} Topological Protection.
        \item \textit{Archetype:} Proton.
    \end{itemize}
    
    \item \textbf{Fundamental Anchors (The Geometric Floor):} Particles defined purely by integer invariants possess \textbf{Infinite Geometric Stability}. They do not decay because they are not "excitations" \textit{of} the lattice; they are the weight of the lattice coordinates themselves. There is no lower-energy geometry to decay into because $D=4$ is a fundamental constant.
\end{enumerate}

\textbf{Validation:}
\begin{itemize}
    \item \textbf{Electron:} $R=0 \le PM$. \textbf{Status:} Geometrically Locked (Stable).
    \item \textbf{Proton:} $R \approx 7.5 \times 10^{-4} > PM$.
    \begin{itemize}
        \item \textit{Analysis:} Geometric fit is imperfect (stress exists), but it is the lightest Baryon ($\Delta^2$). No decay path exists to a lighter $\chi=2$ color singlet.
        \item \textit{Status:} Topologically Protected (Stable).
    \end{itemize}
    \item \textbf{Muon:} $R \approx 9.4 \times 10^{-3} > PM$.
    \begin{itemize}
        \item \textit{Analysis:} Geometric fit is poor. A lighter state exists in the same topological class (Electron, Lepton Number conserved).
        \item \textit{Status:} Unstable. Decays to Ground State.
    \end{itemize}
\end{itemize}

\subsection*{The Decay Rate Function}

For particles that are not locked ($R > PM$), the decay rate $\Gamma$ is determined by the magnitude of the Geometric Residual. We quantify this relationship using a modified Fermi's Golden Rule, where the residual represents the coupling strength to the decay channel.

We identify the specific functional form:
\[ \Gamma \propto \Theta(M - M_{ground}) \cdot |R_{rel}|^\beta \cdot N_{terms} \]

Where:
\begin{itemize}
    \item $\Theta$: The Topological Step Function (0 if Ground State, 1 if Excited).
    \item $|R_{rel}|^\beta$: The \textbf{Geometric Phase Space}. Empirical analysis indicates a power law with \textbf{$\beta \approx 3$}, consistent with 3-body decay phase space volume ($\Gamma \propto Q^5$ where $Q \sim R$).
    \item $N_{terms}$: The \textbf{Entropic Multiplier}. Higher complexity configurations ($N > 2$) possess a larger number of combinatorial decay channels, linearly increasing the total width.
\end{itemize}

\textit{Physical Interpretation:}
\begin{itemize}
    \item \textbf{Residual Dominance:} For fundamental fermions (Low $N$), the lifetime is dominated by the geometric fit ($\tau \propto R^{-3}$).
    \item \textbf{Complexity Dominance:} For high-mass resonances (High $N$, e.g., Top Quark), the complexity multiplier accelerates decay, pushing lifetimes to the yoctosecond scale ($10^{-25}$ s).
\end{itemize}

This framework transforms Particle Decay from a random quantum process into a \textbf{Geometric Error Correction} protocol: the universe works to minimize the residual by cascading high-error states down to the nearest Topological Ground State.

\subsection*{The Eigenstate Selection Rule}

When the configuration space yields multiple geometric formulas for the same mass $M$, the geometry that minimizes Entropic Action ($\lambda \to 0$) will persist the longest.

We define the \textbf{Persistence Index ($P$)} of a configuration as the inverse of its Dissipative Action:
\[ P \propto \frac{1}{N_{terms} \cdot |R_{residual}|} \]

\begin{itemize}
    \item \textbf{The Selection Law:} The physical particle corresponds to the \textbf{Dominant Eigenstate}—the configuration with the highest Persistence Index.
    \item \textbf{Consequence (Tier Promotion):} This rule naturally enforces the generational structure. For the Tau ($M \approx 3477 \, m_e$), the lattice offers a Tier 0 solution ($N=16$) and a Tier 2 solution ($N=3$). Because $N=3 \ll 16$, the Tier 2 configuration has vastly higher persistence. The vacuum "promotes" the resonance to the Areal Tier to minimize the entropic cost of complexity.
\end{itemize}